\section{Matsubara 格林函数：谱表示与频率求和}

有限温度计算的标准语言是 Matsubara 技巧（Mahan 第 3 章）。本节给出周期性、谱表示、以及最常用的频率求和公式——这是图计算真正动手时的核心代数。

\subsection{虚时定义与（反）周期性}

巨正则密度矩阵 $\rho=e^{-\beta(H-\mu N)}/Z$。定义
\begin{equation}
  \mathcal{G}(\tau)
  =-\bigl\langle T_\tau\,c(\tau)c^\dagger(0)\bigr\rangle,
  \qquad
  c(\tau)=e^{\tau K}c\,e^{-\tau K},
  \quad K=H-\mu N,
\end{equation}
$\tau\in[0,\beta]$。利用迹的循环性：
\begin{equation}
\begin{aligned}
  \mathcal{G}(\tau<0)
  &=+\bigl\langle c^\dagger(0)c(\tau)\bigr\rangle
  =+\frac{1}{Z}\mathrm{Tr}\bigl[e^{-\beta K}c^\dagger e^{\tau K}c e^{-\tau K}\bigr]\\
  &=-\zeta\,\mathcal{G}(\tau+\beta),
\end{aligned}
\end{equation}
$\zeta=-1$ 费米、$+1$ 玻色。故费米子\textbf{反周期}、玻色子周期：
\begin{equation}
  \mathcal{G}(\tau+\beta)=\zeta\,\mathcal{G}(\tau).
\label{eq:Gperiod}
\end{equation}

\subsection{Matsubara 频率}

Fourier 级数
\begin{equation}
  \mathcal{G}(\tau)
  =\frac{1}{\beta}\sum_n
  e^{-\mathrm{i}\omega_n\tau}\mathcal{G}(\mathrm{i}\omega_n),
  \qquad
  \mathcal{G}(\mathrm{i}\omega_n)
  =\int_0^\beta\mathrm{d}\tau\,
  e^{\mathrm{i}\omega_n\tau}\mathcal{G}(\tau),
\end{equation}
\begin{equation}
  \omega_n
  =
  \begin{cases}
    (2n+1)\pi/\beta, & \text{费米},\\
    2n\pi/\beta, & \text{玻色},
  \end{cases}
  \qquad n\in\mathbb{Z}.
\end{equation}

\subsection{谱表示}

插入 $K$ 的本征态，与零温 Lehmann 同样步骤，得到
\begin{equation}
  \mathcal{G}(\mathrm{i}\omega_n)
  =\int_{-\infty}^{\infty}
  \frac{\mathrm{d}\omega\,A(\omega)}{\mathrm{i}\omega_n-\omega},
\label{eq:G_spectral}
\end{equation}
$A(\omega)\ge0$，$\int A=1$。解析延拓的唯一性（在无穷远处 $G\sim 1/z$）保证
\begin{equation}
  G^R(\omega)
  =\mathcal{G}(\mathrm{i}\omega_n\to\omega+\mathrm{i}0^+)
  =\int\frac{A(\omega')}{\omega-\omega'+\mathrm{i}0^+}\mathrm{d}\omega'.
\end{equation}
自由费米子 $A(\omega)=\delta(\omega-\xi_{\mathbf{k}})$，$\xi_{\mathbf{k}}=\varepsilon_{\mathbf{k}}-\mu$，
\begin{equation}
  \mathcal{G}_0(\mathbf{k},\mathrm{i}\omega_n)
  =\frac{1}{\mathrm{i}\omega_n-\xi_{\mathbf{k}}},
  \qquad
  \mathcal{G}_0(\mathbf{k},\tau)
  =-e^{-\xi_{\mathbf{k}}\tau}
  \bigl[\theta(\tau)(1-n_F(\xi_{\mathbf{k}}))
  -\theta(-\tau)n_F(\xi_{\mathbf{k}})\bigr].
\label{eq:G0_tau}
\end{equation}

\subsection{频率求和：留数技巧}

图中出现
\begin{equation}
  S=\frac{1}{\beta}\sum_n
  h(\mathrm{i}\omega_n).
\end{equation}
费米情形用 $n_F(z)=(e^{\beta z}+1)^{-1}$，其极点在 $z=\mathrm{i}\omega_n$，留数 $-1/\beta$。故
\begin{equation}
  \frac{1}{\beta}\sum_n h(\mathrm{i}\omega_n)
  =-\sum_j\mathrm{Res}\bigl[h(z)n_F(z)\bigr]_{z=z_j},
\label{eq:matsum_F}
\end{equation}
$z_j$ 为 $h$ 的极点（围道在无穷远处收敛需 $h(z)\sim 1/z^{1+\delta}$；否则乘 $e^{z0^+}$）。玻色求和把 $n_F$ 换成 $-n_B(z)$，$n_B=(e^{\beta z}-1)^{-1}$。

\begin{example}[占据数]
取 $h(z)=e^{z\eta}/(z-\xi)$，$\eta\to0^+$，
\begin{equation}
  \frac{1}{\beta}\sum_n
  \frac{e^{\mathrm{i}\omega_n\eta}}{\mathrm{i}\omega_n-\xi}
  =n_F(\xi).
\end{equation}
这正是 $\langle c^\dagger c\rangle=n_F(\xi)$。$\eta\to0^-$ 则给出 $n_F-1$。
\end{example}

\begin{example}[极化气泡中的费米圈]
RPA/Lindhard 的 Matsubara 形式
\begin{equation}
  \Pi(\mathbf{q},\mathrm{i}\nu_m)
  =\frac{2}{\beta L^d}
  \sum_{\mathbf{k},n}
  \mathcal{G}_0(\mathbf{k},\mathrm{i}\omega_n)
  \mathcal{G}_0(\mathbf{k}+\mathbf{q},\mathrm{i}\omega_n+\mathrm{i}\nu_m).
\end{equation}
对 $n$ 用 \eqref{eq:matsum_F}，
\begin{equation}
  \Pi(\mathbf{q},\mathrm{i}\nu_m)
  =\frac{2}{L^d}
  \sum_{\mathbf{k}}
  \frac{n_F(\xi_{\mathbf{k}})-n_F(\xi_{\mathbf{k}+\mathbf{q}})}
  {\mathrm{i}\nu_m+\xi_{\mathbf{k}}-\xi_{\mathbf{k}+\mathbf{q}}},
\label{eq:Pi_matsum}
\end{equation}
解析延拓后即 Lindhard 函数 \eqref{eq:lindhard_def}。这是 Mahan 把图规则接到电子气响应的标准步骤。
\end{example}

\subsection{自由玻色（声子）传播子}

对位移型算符 $A_{\mathbf{q}}=a_{\mathbf{q}}+a_{-\mathbf{q}}^\dagger$，
\begin{equation}
  \mathcal{D}_0(\mathbf{q},\mathrm{i}\nu_m)
  =-\frac{2\omega_{\mathbf{q}}}{\nu_m^2+\omega_{\mathbf{q}}^2}
  =\frac{1}{\mathrm{i}\nu_m-\omega_{\mathbf{q}}}
  -\frac{1}{\mathrm{i}\nu_m+\omega_{\mathbf{q}}}.
\label{eq:D0}
\end{equation}
（总体符号随 $\mathcal{D}=-\langle T_\tau A A\rangle$ 的定义而异；上式与常见电子--声子自能公式配套。）
